% Figure: a contact network with heterogeneous degree, and the contagion
% curve on the network compared with the homogeneous-mixing mean-field SIR.
\begin{figure}[htbp]
\centering
\begin{tikzpicture}[font=\small]
% ===== Left: the contact network =====
\begin{scope}[xshift=0cm]
\node[font=\small, anchor=south west] at (-0.3, 4.0)
  {\textbf{(a) a contact network}};
% a small heterogeneous graph: one hub plus periphery
\tikzstyle{node0}=[circle, draw=engblack!70, fill=engblue!12,
  minimum size=4mm, inner sep=0pt]
\tikzstyle{hub}=[circle, draw=engblack!80, fill=engred!30,
  minimum size=7mm, inner sep=0pt, font=\scriptsize\bfseries]
\node[hub]   (h) at (2.0,2.0) {hub};
\node[node0] (a) at (0.2,3.4) {};
\node[node0] (b) at (0.0,2.0) {};
\node[node0] (c) at (0.4,0.7) {};
\node[node0] (d) at (1.8,0.2) {};
\node[node0] (e) at (3.4,0.6) {};
\node[node0] (f) at (3.9,2.1) {};
\node[node0] (g) at (3.3,3.5) {};
\node[node0] (i) at (1.9,3.8) {};
\node[node0] (j) at (4.4,3.4) {};
\foreach \n in {a,b,c,d,e,f,g,i} { \draw[draw=engblack!45] (h) -- (\n); }
\draw[draw=engblack!45] (a) -- (b);
\draw[draw=engblack!45] (b) -- (c);
\draw[draw=engblack!45] (f) -- (j);
\draw[draw=engblack!45] (g) -- (j);
\draw[draw=engblack!45] (f) -- (g);
\node[font=\scriptsize, anchor=north] at (2.0,-0.3)
  {degree-heterogeneous: a few high-degree hubs};
\end{scope}
% ===== Right: contagion curve, network vs mean-field =====
\begin{scope}[xshift=7.2cm]
\node[font=\small, anchor=south west] at (-0.3, 4.0)
  {\textbf{(b) infectious fraction over time}};
\draw[->, thick, draw=engblack!70] (0,0) -- (6.0,0)
  node[right, font=\scriptsize] {time $t$};
\draw[->, thick, draw=engblack!70] (0,0) -- (0,3.6)
  node[above, font=\scriptsize] {$I(t)/N$};
% mean-field SIR: later, taller peak
\draw[very thick, draw=engblue!80, smooth, samples=80, domain=0:5.6]
  plot (\x, { 3.0*0.46*exp(-((\x-3.0)*(\x-3.0))/0.9) });
% network: earlier, lower, fatter peak (hub seeds faster, heterogeneity
% blunts the peak)
\draw[very thick, draw=engred!85, dashed, smooth, samples=80, domain=0:5.6]
  plot (\x, { 3.0*0.30*exp(-((\x-2.2)*(\x-2.2))/1.7) });
\draw[very thick, draw=engblue!80] (3.4,3.3) -- (3.8,3.3);
\node[anchor=west, font=\scriptsize] at (3.85,3.3) {mean-field};
\draw[very thick, draw=engred!85, dashed] (3.4,2.9) -- (3.8,2.9);
\node[anchor=west, font=\scriptsize] at (3.85,2.9) {network sim.};
\end{scope}
\end{tikzpicture}
\caption[A contact network and the contagion it carries]{(a) A contact
network in which most individuals have low degree and a few hubs have
high degree, the degree heterogeneity that homogeneous mixing erases. The
infection threshold on such a network is governed not by the mean degree
but by the ratio of the second moment of the degree distribution to the
first, which the hubs inflate. (b) The infectious fraction over time,
computed on the network by individual-level stochastic simulation
(dashed) and by the homogeneous-mixing mean-field SIR of
section~\ref{fig:vol02:ch18:compartment-sir} (solid). The hubs seed the
epidemic earlier, so the network curve rises sooner; the degree
heterogeneity spreads the infection across a wider range of effective
contact rates, so the network peak is lower and broader than the
mean-field model predicts. The gap between the two curves is the
structural error the homogeneous-mixing assumption introduces, and its
sign and shape are the network model's specification.}
\label{fig:vol02:ch18:network-contagion}
\end{figure}
